\documentclass[11pt,a4paper]{article}
\usepackage[T1]{fontenc}
\usepackage[utf8]{inputenc}
\usepackage{amsmath,amssymb,xcolor}
\usepackage[margin=25mm]{geometry}
\usepackage[colorlinks=true,urlcolor=blue]{hyperref}
\definecolor{nexorablue}{HTML}{2459B5}
\setlength{\parindent}{0pt}
\setlength{\parskip}{7pt}
\setlength{\emergencystretch}{2em}
\renewcommand{\familydefault}{\sfdefault}
\begin{document}
\makeatletter
\def\@oddfoot{\footnotesize Download for free at \url{https://openstax.org/details/books/university-physics-volume-2}.\hfil\thepage}
\let\@evenfoot\@oddfoot
\makeatother
{\small\bfseries\color{nexorablue} NEXORA CPGE}\par
{\small MPSI / OST-SUP-10-09 / Énoncé}\par
{\LARGE\bfseries Régler l’intervalle des essuie-glaces}\par
\vspace{0.25em}\hrule\vspace{0.6em}

Une source idéale de $12\,\mathrm V$ charge un condensateur de $10{,}00\,\mathrm{mF}$ à travers un rhéostat réglable de $0$ à $10{,}00\,\mathrm{k\Omega}$. Lorsque la tension du condensateur atteint $10{,}00\,\mathrm V$, un interrupteur ferme le circuit du moteur et ramène la tension à zéro. Négliger la durée de cette remise à zéro. Quel réglage impose un balayage toutes les $10{,}00\,\mathrm s$ ?

\bigskip
{\footnotesize Adaptation française par Nexora CPGE. Source : Samuel J. Ling, Jeff Sanny et William Moebs (OpenStax, Rice University),
\href{https://d3bxy9euw4e147.cloudfront.net/oscms-prod/media/documents/UniversityPhysicsVolume2-OP.pdf\#page=484}{\textit{University Physics Volume 2, édition 2018}}, § 10.5, exemple 10.9, p. 474.
\href{https://creativecommons.org/licenses/by/4.0/}{CC BY 4.0}.}
\end{document}
